\documentclass[10pt,a4paper]{amsart}
\usepackage[margin=19mm]{geometry}
\usepackage{amsmath,amssymb,mathtools}
\usepackage[scr=rsfs,cal=euler]{mathalfa}
\usepackage{xfrac}
\usepackage[unicode,hidelinks,bookmarks=false]{hyperref}
\newcommand{\ie}{i.e.\ }
\newcommand{\cf}{cf.\ }
\newcommand{\resp}{respectively\ }
\newcommand{\Subsection}{\subsection}
\providecommand{\nobreakdash}{\nobreak\hbox{-}}
\setlength{\emergencystretch}{2em}
\multlinegap0pt
\usepackage{bm}
\usepackage{bbm}
\usepackage{dsfont}
\usepackage{xspace}
\usepackage{cancel}
\usepackage{mathtools}
\usepackage{amsthm}
\makeatletter
\@ifundefined{thm}{\newtheorem{thm}{Thm}}{}
\@ifundefined{lemma}{\newtheorem{lemma}[thm]{Lemma}}{}
\makeatother
\newcommand{\be}{\begin{equation}}
\newcommand{\ee}{\end{equation}}
\newcommand{\rmd}{{\rm d}}
\title[On the existence of fibered three-dimensional perfect fluid ]{On the existence of fibered three-dimensional perfect fluid equilibria without continuous Euclidean symmetry}
\author{Theodore D. Drivas, Tarek M. Elgindi, Daniel Ginsberg}
\date{Source: arXiv:2510.02955v2 (CC BY 4.0)}
\begin{document}
\maketitle

\noindent\emph{Source and licence.} On the existence of fibered three-dimensional perfect fluid equilibria without continuous Euclidean symmetry. Version arXiv:2510.02955v2 (CC BY 4.0), distributed under the Creative Commons Attribution 4.0 International licence, as recorded in the arXiv licence metadata for this exact version and shown on the abstract page https://arxiv.org/abs/2510.02955v2. Full rights notice: licenses/arxiv-2510.02955-CC-BY-4.0.txt.\par\smallskip

\noindent\textit{Editorial scope.} Counted unit: the Lemma whose source label is \texttt{Xpoin} in the article (source0.tex lines 547--566, source label \texttt{Xpoin}), delivered together with the article's own complete original proof of it (source0.tex lines 567--609, one \texttt{proof} environment of 43 lines). No mathematics is written, repaired or re-derived: statement and proof are the article's own text line for line. Required context retained uncounted: none. Every definition, display and label of those lines is retained unchanged. Omitted from this package and not redistributed: 1--546, 610--820. Nothing inside a retained excerpt is omitted or altered: every retained body is delivered exactly as the article's lines are written, so no omission receipt of any kind accompanies this package. Citations: the retained excerpts carry no citation command at all, so no bibliography entry of the article is transcribed and no reference is redistributed. Reference adaptation: none was needed and none was made; every reference inside the delivered unit resolves inside it, so no unresolved reference or citation is produced. Printed slips kept literal: no printed word, symbol, hypothesis or number of the counted statement or of its proof was altered. Layout and class: the article's own class is \texttt{article} with packages including amsmath, amssymb, amsthm, amsfonts, color, longtable, hyperref, enumitem, setspace, babel, csquotes, graphicx, latexsym, array; the wrapper is the lane's pinned \texttt{amsart} editorial wrapper at 10pt on A4, which restates the theorem-like environments the retained text needs and adds the article's own macro definitions that the retained text needs (\texttt{\textbackslash be}, \texttt{\textbackslash ee}, \texttt{\textbackslash rmd}), with extra wrapper packages (bm, bbm, dsfont, xspace, cancel, mathtools). The delivered numbering is the unit's own. Assets: no retained span contains a figure environment, a graphics-inclusion command, a PDF-page inclusion, a table or a TikZ picture; no image file is copied into this package, the package declares no asset dependency and no image file is redistributed with it. Licence: version arXiv:2510.02955v2 of this article is distributed under the Creative Commons Attribution 4.0 International licence, as recorded in the arXiv licence metadata for this exact version and shown on the abstract page https://arxiv.org/abs/2510.02955v2. A full rights notice is delivered at licenses/arxiv-2510.02955-CC-BY-4.0.txt. Characters: every retained excerpt is pure ASCII, so the delivered engine, XeLaTeX with its default Latin Modern text font, reports no missing character.
\begin{lemma}
  \label{Xpoin}
  Let $U$ be a domain foliated by a family of simple, $C^{k,\alpha}$-smooth
  and $m-$fold symmetric curves; that is, for each curve
  $\gamma$, there is a curve $\gamma_0$ so that
  $\gamma = \bigcup_{j=0}^{m-1}  O_{2\pi/m}^j\gamma_0$.
  Let $X$ be a $C^{k,\alpha}$ vector field on $U$ so that $X|_\gamma$ is tangent
  to $\gamma$ for each curve $\gamma$ in the foliation, and that the $X$ period $T(\gamma) := \int_\gamma \frac{\rmd s}{|X|}\leq C_{\mathsf{per}}$ is bounded uniformly.
  Further assume that $X$ is $m$-fold symmetric in the sense that $X= ({O_{2\pi/m}})_* X$
  where $({O_{2\pi/m}})_* $ denotes the pushforward.  Then, if $f$ is an $m$-fold symmetric function with $\int_{\gamma} f\, \rmd s = 0$ for each
  curve $\gamma$ in the foliation, then there is a constant
  $C = C\left(C_{\mathsf{per}},\|X\|_{C^{k,\alpha}}, \|\gamma\|_{C^{k,\alpha}},
     k, \alpha\right)$ so 

  \begin{equation}
    \label{poin}
    \| f \|_{C^{k,\alpha}(U)} \leq   \frac{C}{m} \| \nabla_X f\|_{C^{k,\alpha}(U)}.
  \end{equation}

\end{lemma}

\begin{proof}

    By the $m$-fold symmetry, we can write $U = \bigcup_{j = 0}^{m-1} (O_{2\pi/m})^j U_0$
    for a domain $U_0$, and it suffices
    to prove the bound \eqref{poin} with $U$ replaced by $U_0$. Fix a curve
    $\gamma$ in the foliation and let $\gamma_0$ denote its restriction to $U_0$. Let $x,y\in \gamma_0$ and let $\Phi_t$ be the flowmap for $X$, i.e.
$$
\frac{\rmd }{\rmd t} \Phi_t = X\circ \Phi_t , \qquad \Phi_0={\rm id}.
$$
Let $T(x,y)$ be such that $\Phi_{T(x,y)}(x) = y$. 
Since $f$ is mean zero on each $\gamma$ and is $m$-fold symmetric,
    $f$ is mean zero on $\gamma_0$, so
integrating the expression $f(x)- f(y) =\int_0^{T(x,y)} (\nabla_X  f)\circ \Phi_t(x) d t$
in $y$ over $\gamma_0$, we find
\begin{equation} \label{finvform}
f(x) = \frac{1}{{\rm length}(\gamma_0)}\int_{\gamma_0}\int_0^{T(x,y)} (\nabla_X  f)\circ \Phi_t(x) \rmd t \rmd y,
\end{equation}
and it follows that for any $x$ on $\gamma_0$,
\begin{align*}
|f(x)| \leq \frac{1}{{\rm length}(\gamma_0)}\int_{\gamma_0}\int_0^{T(x,y)} |(\nabla_X  f)\circ \Phi_t(x)| \rmd t \rmd y 
%&=  \frac{1}{{\rm length}(\gamma_0)}\int_{\gamma_0}\int_{\gamma_0} |\nabla_X  f| \frac{ds}{| X|} dy\\
&\leq  T(\gamma_0)\| \nabla_X f\|_{L^\infty}
\end{align*}
where 
\be
T(\gamma_0):=  \int_{\gamma_0} \frac{\rmd s}{|X|}
\ee
is the time taken to traverse the segment $\gamma_0$, end-to-end. By the $m$-fold
symmetry, $T_*(\gamma_0)
= \frac{1}{m}T_*(\gamma)$. Since the curves $\gamma_0$
foliate the domain $U_0$, this gives the bound when
$k = \alpha = 0$. We now show how to bound the first derivative, higher-order derivatives
being similar.  We express the line integral via its parametrization by the flow of $X$, namely $\Phi_t(x):[s_1(x),s_2(x)] \to \gamma_0$:
\begin{align}
f(x) &= \frac{1}{\int_{s_1(x)}^{s_2(x)} |X\circ \Phi_s(x)|  ds}\int_{s_1(x)}^{s_2(x)} \int_0^{T(x,\Phi_s(x))} (\nabla_X  f)\circ \Phi_t(x) |X\circ \Phi_s(x)| \rmd t \rmd s,
\end{align}
Note that derivatives of $s_i(x)$ and $\Phi_t(x)$ are bounded by derivatives of $X$.  Thus
\begin{align*}
|\nabla f(x)| &\leq  \Bigg[\frac{1}{m} T(\gamma) \|s_i\|_{C^1} \|X\|_{L^\infty} \|\nabla_X  f\|_{L^\infty}  +\frac{\sup_{t\in [s_1,s_2]}\|\nabla_x T(\cdot,\Phi_t (\cdot) )\|_{L^\infty}}{{\rm length}(\gamma_0)} \int_{\gamma_0}|\nabla_X  f(y)| d y \\
&\qquad + \frac{\sup_t\| \nabla \Phi_t(\cdot ) \|_{L^\infty}}{{\rm length}(\gamma_0)}\int_{\gamma_0}\int_0^{T(x,y)} |\nabla \nabla_X  f\circ \Phi_t(x)| \rmd t \rmd y\Bigg].
\end{align*}
Note that, if $f$ is  $m$-fold symmetric, then so is $\nabla_X f$ and, moreover, it is mean zero on each $\gamma_0$.  As such, can can apply the estimate for $\|f\|_{L^\infty}$ to arrive at $\|\nabla_X f\|_{L^\infty}\lesssim 1/m \|\nabla \nabla_X f\|_{L^\infty}.$ Likewise, the last term is bounded by the same argument.  This concludes the sketch of the argument.  Higher derivatives follow similarly.
\end{proof}

\end{document}
